\section[Power Multipliers]{Power Multipliers and Radial Perturbations}
\label{app:power-identities}

All identities involving negative or nonintegral powers are used on
\(R,S>0\).  In the power calculations, \(F\) is homogeneous of
degree \(e\geq1\) and \(\cL(F)=F\lambda_F\), with \(\lambda_F\)
smooth on the calculation domain.  It is a polynomial whenever the
divisibility hypothesis of \cref{thm:exact-multiplier} holds.

\subsection{Notation and power factors}

Consider the power family
\begin{equation}\label{eq:power-family}
 G_{\alpha,\gamma}=R^\alpha S^\gamma F,
 \qquad R=|x|^2,\quad S=|\nabla F|^2,
\end{equation}
with real exponents.
Write \(e=\deg F\), and use the contractions
\[
 U=\Gamma(F,S),\quad V=\Gamma(S,S),\quad
 D=\Delta F,\quad B=\Delta S.
\]
Here \(U,V,D,B\) are ambient functions; in a finite Laplacian closure
they are functions of the generators.  Define
\begin{align}
 \mathcal V_\gamma
 &=S+\frac{2\gamma FU}{S}+\frac{\gamma^2F^2V}{S^2},
 \label{eq:gradient-power-V}\\
 \mathcal D_\gamma
 &=D+\frac{2\gamma U}{S}
       +\frac{\gamma FB}{S}
       +\frac{\gamma(\gamma-1)FV}{S^2}.
 \label{eq:gradient-power-D}
\end{align}
Let \(\mathcal Q_\gamma\) be the explicit expression in
\eqref{eq:gradient-power-Q}.  It satisfies
\(\cL(S^\gamma F)=S^{3\gamma}F\mathcal Q_\gamma\).

\subsection{Gradient and radial identities}

\begin{proposition}[The gradient-power defect]
\label{prop:gradient-power-identity}
For \(H=S^\gamma F\),
\[
 |\nabla H|^2=S^{2\gamma}\mathcal V_\gamma,\qquad
 \Delta H=S^\gamma\mathcal D_\gamma,\qquad
 \cL(H)=S^{3\gamma}F\mathcal Q_\gamma,
\]
where \(\mathcal V_\gamma,\mathcal D_\gamma\) are given in
\eqref{eq:gradient-power-V}--\eqref{eq:gradient-power-D}, and
\begin{align}
 \mathcal Q_\gamma={}&\lambda_F+\gamma B
 +\frac{2\gamma UD-\gamma\Gamma(F,U)
                 -(\gamma^2+\tfrac32\gamma)V}{S}
 +\frac{\gamma(\gamma+1)U^2}{S^2}\notag\\
 &+\frac{\gamma^2F}{S^2}
       \left(2UB+VD-\Gamma(S,U)-\tfrac12\Gamma(F,V)\right)\notag\\
 &+\frac{\gamma^3F^2}{S^3}
       \left(VB-\tfrac12\Gamma(S,V)\right).
 \label{eq:gradient-power-Q}
\end{align}
At every regular zero of \(F\), one has \(S>0\). Under the assumed
smoothness of \(\lambda_F\), the displayed formula is therefore
well-defined there; no division by \(F\) remains.
\end{proposition}

\begin{proof}
The gradient and Laplacian follow from the product rule.
For the defect, insert \(a=S^\gamma\) in
\cref{thm:exact-multiplier}.  The derivative identities are
\[
 \nabla a=\gamma S^{\gamma-1}\nabla S,\qquad
 \Hess a=\gamma S^{\gamma-1}\Hess S
       +\gamma(\gamma-1)S^{\gamma-2}\,dS\otimes dS.
\]
Differentiating \(S=|\nabla F|^2\), \(U=\Gamma(F,S)\), and
\(V=|\nabla S|^2\) gives
\begin{align*}
 \Hess F(\nabla F,\nabla S)&=V/2,\\
 \Hess S(\nabla F,\nabla F)&=\Gamma(F,U)-V/2,\\
 \Hess F(\nabla S,\nabla S)&=\Gamma(S,U)-\Gamma(F,V)/2,\\
 \Hess S(\nabla F,\nabla S)&=\Gamma(F,V)/2,\\
 \Hess S(\nabla S,\nabla S)&=\Gamma(S,V)/2.
\end{align*}
Substitution and collection of the powers of \(F\) yield
\eqref{eq:gradient-power-Q}.
\end{proof}

\begin{proposition}[The radial master identity]
\label{prop:radial-master-identity}
Let \(H\) be smooth and homogeneous of degree \(m\), and set
\(S_H=|\nabla H|^2\), \(D_H=\Delta H\), and
\(c_\alpha=4\alpha(m+\alpha)\).  Then
\begin{align}
 |\nabla(R^\alpha H)|^2
 &=R^{2\alpha}\left(S_H+\frac{c_\alpha H^2}{R}\right),
 \label{eq:radial-master-gradient}\\
 \cL(R^\alpha H)
 &=R^{3\alpha}\Biggl[\cL(H)
  +\frac{c_\alpha H^2D_H}{R}
  +\frac{\bigl(2\alpha(N-1)-c_\alpha\bigr)HS_H}{R}\notag\\
 &\hspace{31mm}
  +\frac{c_\alpha\bigl(m+2\alpha(N-1)\bigr)H^3}{R^2}
  \Biggr].
 \label{eq:radial-master-L}
\end{align}
\end{proposition}

\begin{proof}
Euler's identity gives \(\Gamma(R,H)=2mH\).  Hence
\[
 \Delta(R^\alpha H)
 =R^\alpha\left[D_H+
    \frac{2\alpha(2m+N+2\alpha-2)H}{R}\right],
\]
and the product rule gives \eqref{eq:radial-master-gradient}.
Writing \(W=S_H+c_\alpha H^2/R\), of degree \(2m-2\), yields
\[
 R^{-3\alpha}\cL(R^\alpha H)
 =WD_H+\frac{2\alpha(N-1)HW}{R}-\tfrac12\Gamma(H,W).
\]
Use
\(\Gamma(H,H^2/R)=2HS_H/R-2mH^3/R^2\)
to expand the last term.  This proves \eqref{eq:radial-master-L}.
\end{proof}



\subsection{The combined multiplier}

\begin{proposition}[The combined power multiplier]
\label{prop:combined-power-multiplier}
Suppose \(\cL(F)=F\lambda_F\).  Set
\[
 m=e+2\gamma(e-1),\quad \ell=m+2\alpha,\quad
 c_\alpha=4\alpha(m+\alpha),\quad
 b_\alpha=2\alpha(N-1)-c_\alpha.
\]
For \(G=G_{\alpha,\gamma}\), one has
\begin{align}
 |\nabla G|^2
 &=R^{2\alpha}S^{2\gamma}\widehat V_{\alpha,\gamma},
 &\widehat V_{\alpha,\gamma}
 &=\mathcal V_\gamma+\frac{c_\alpha F^2}{R},
 \label{eq:combined-power-V}\\
 \cL(G)
 &=R^{3\alpha}S^{3\gamma}F\widehat Q_{\alpha,\gamma},
 \label{eq:combined-power-L}
\end{align}
where
\begin{equation}\label{eq:combined-power-Q}
 \widehat Q_{\alpha,\gamma}
 =\mathcal Q_\gamma+\frac{c_\alpha F\mathcal D_\gamma}{R}
 +\frac{b_\alpha\mathcal V_\gamma}{R}
 +\frac{c_\alpha\bigl(m+2\alpha(N-1)\bigr)F^2}{R^2}.
\end{equation}
Thus \(\widehat V_{\alpha,\gamma}>0\) and
\(\widehat Q_{\alpha,\gamma}\geq0\) give the normalized-gradient
sign on the region where they hold.  If that region is one phase,
they give its one-sided sign.  Comparison across the omitted set
requires the hypotheses on the exceptional set and boundary flux in
\cref{thm:paired-subcalibration}.
\end{proposition}

\begin{proof}
Apply \cref{prop:radial-master-identity} to \(H=S^\gamma F\), of
degree \(m\), using the three identities of
\cref{prop:gradient-power-identity}. This gives
\eqref{eq:combined-power-V}--\eqref{eq:combined-power-Q}.
Since \(R^\alpha S^\gamma>0\), \(G\cL(G)\) has the sign of
\(F^2\widehat Q_{\alpha,\gamma}\).
\end{proof}

\subsection{Radial perturbations giving strict divergence estimates}

\begin{proposition}[Moving the degree toward \(n/2\)]
\label{prop:radial-strictification}
Let \(G\in C^\infty(\R^{n+1}\setminus\{0\})\) be homogeneous of degree
\(\ell\neq n/2\), with \(\nabla G\neq0\) throughout the punctured space.
Suppose \(C=\{G=0\}\) has a smooth compact minimal link,
\(G\) takes both signs, and
\[
 \operatorname{sgn}(G)\diver X_G>0\quad\text{off }C,
 \qquad\beta_G\geq0\quad\text{on }C\cap\Sph^n,
\]
where \(\beta_G\) is defined in \eqref{eq:boundary-Jacobi-potential}.
Put \(\sigma=\operatorname{sgn}(\ell-n/2)\).  For all sufficiently small
\(\eps>0\), the potential
\begin{equation}\label{eq:strictifying-potential}
 G_\eps=r^{-\sigma\eps}G,\qquad \ell_\eps=\ell-\sigma\eps,
\end{equation}
has nonvanishing gradient on the punctured space and generates a global
subcalibration satisfying
\begin{equation}\label{eq:strictified-distance-bound}
 \operatorname{sgn}(G_\eps)\diver X_{G_\eps}
 \geq c_\eps\frac{\dist(x,C)}{|x|^2}
 \qquad(x\neq0)
\end{equation}
for a constant \(c_\eps>0\).  Writing \(\beta_\eps=\beta_{G_\eps}\), its boundary coefficient is
\begin{equation}\label{eq:strictified-boundary}
 \beta_\eps=\beta_G+\left(\ell-\frac n2\right)^2
                 -\left(\ell_\eps-\frac n2\right)^2
       =\beta_G+2\left|\ell-\frac n2\right|\eps-\eps^2.
\end{equation}
The new potential is generally not polynomial. If \(G\) is polynomial
and \(\eps\) is positive rational, then \(G_\eps\) is semialgebraic
on \(r>0\): the radial factor is a rational power of the positive
polynomial \(R\), hence is described by a polynomial equation and
the choice of its positive root.
\end{proposition}

\begin{proof}
Nonvanishing of \(\nabla G_\eps\) for small \(\eps\) follows from smooth
parameter dependence and compactness of the unit sphere.
On \(C\), its reciprocal gradient length is
\(w_\eps=r^{\sigma\eps}w=r^{1-\ell_\eps}\phi\).
For \(w=r^\alpha\phi\), separation of variables gives
\[
 -r^2\frac{J_Cw}{w}
 =\frac{(-\Delta_\Sigma-|A_\Sigma|^2)\phi}{\phi}
                  -\alpha(\alpha+n-2).
\]
Use \(\alpha=1-\ell\) and \(\alpha_\eps=1-\ell_\eps\) to obtain
\eqref{eq:strictified-boundary}.

For the global sign, let \(s\) be signed spherical distance to the link
in a fixed tubular neighborhood, positive on \(G>0\).  On \(r=1\), put
\(f_\eps=\diver X_{G_\eps}\).
Minimality gives \(f_\eps=0\) on the link.  The quotient
\(h_\eps=f_\eps/s\) extends smoothly across it, with
\(h_\eps|_\Sigma=\beta_\eps\).
By \eqref{eq:strictified-boundary},
\(\partial_\eps h_\eps|_{\eps=0,\Sigma}=2|\ell-n/2|>0\).
Compactness and continuity give a smaller fixed tube and a parameter
interval on which \(\partial_\eps h_\eps\geq|\ell-n/2|\).
Since \(h_0\geq0\), integration in the parameter yields
\(h_\eps\geq\eps|\ell-n/2|\) throughout that tube.
On the complement of a still smaller tube the original strict sign has a
positive minimum, which persists for small \(\eps\).
Spherical normal distance is comparable to \(\dist(\theta,C)\) near the
link; away from it the latter is bounded.  This proves the uniform
unit-sphere estimate, and homogeneity gives
\eqref{eq:strictified-distance-bound}.
\end{proof}

For a singular link,
\eqref{eq:strictified-boundary} still holds on the regular part.
The global sign requires additional estimates near the singular strata.

\subsection{Radial multiplication preserving the divergence sign}

\begin{proposition}[Radial multiplication of a homogeneous potential]
\label{prop:cubic-radial-lift}
Let \(h\) be smooth and homogeneous of degree \(d>0\) on a conical
open subset of \(\R^N\setminus\{0\}\). If
\[
 h\cL(h)\ge0,\qquad h\Delta h\ge0,\qquad N\ge2d+2,
\]
then \(g=\sqrt R\,h\) satisfies \(g\cL(g)\ge0\), strictly where
\(h\ne0\). If \(\nabla h\ne0\) throughout the domain, then
\(\nabla g\ne0\) there as well.
\end{proposition}

\begin{proof}
The radial master identity, \cref{prop:radial-master-identity} with
\(\alpha=1/2\), reads
\begin{align}
 \cL(\sqrt R\,h)=R^{3/2}\Biggl[\cL(h)
 &+(2d+1)\frac{h^2\Delta h}{R}
 +(N-2d-2)\frac{h|\nabla h|^2}{R}\notag\\
 &+(2d+1)(d+N-1)\frac{h^3}{R^2}\Biggr].
 \label{eq:cubic-radial-lift}
\end{align}
Multiplication by \(g=\sqrt R\,h\) makes every term nonnegative;
the last one is positive when \(h\ne0\). The same master identity
gives
\[
 |\nabla g|^2=R|\nabla h|^2+(2d+1)h^2,
\]
proving the final assertion.
\end{proof}

